% v16: proof-derived finite reduction for the non-B representation of leaf 4.
% This file is intended to follow the global z=1, r<=4 reduction.

\subsection{単位 entry の枝の恒等式と有限化}
\label{zu16:section}

本節では交差標識 $(A,B,C,D)=(0,1,4,3)$ を固定し，
$z=1$，$1\le r\le4$，$S>r$，$rT>(r+1)S$ を仮定する．
前節までの局所 chart，整数 clutching，node 非負性，対向係数の相殺を保持する．
葉 $4$ の ray $2$ を一高分母標準形の $C$ または $D$ として表せる場合を扱う．
分母三の標識重複によって $B$ としても表せる行列を除外するという意味ではない．

\begin{lemma}[単位 entry の二つの恒等式]
\label{zu16:identities}
$T=(r+1)y$，$x=sS/r$ と置くと，$y,x$ は正整数であり，
\begin{align}
 (r+2)sS-r(rs+1)y&=2r(s-1),\label{zu16:pair}\\
 F_r(s)S&=r(rs+1)E+2r(r+2)(s-1),\label{zu16:E}\\
 F_r(s)&=-rs^2+(4r+3)s-r,\qquad
 E=D^2_4=(r+2)y-\frac{s+r}{r}S.\notag
\end{align}
$1\le r\le4$，$s\in\ZZ_{>0}$ では $F_r(s)\ne0$ である．
さらに $E\ge-2$ なら $s\le9$ である．
\end{lemma}

\begin{proof}
$A=a+x-s$，$B=b-y+1$，$a=sS+1$，$b=(r+1)y+1$ を
$AB=ab-1$ に代入して $r$ を掛けると式~\eqref{zu16:pair} を得る．
それと $E=(r+2)y-(s+r)S/r$ から $y$ を消去すると
式~\eqref{zu16:E} が従う．
$s=1,2,3,4$ での $F_r(s)$ は順に
$2r+3,3r+6,2r+9,12-r$ で正である．
$s=5,6$ では $15-6r,18-13r$ であり，$r=1,2,3,4$ のどれでも零でない．
$s\ge7$ では $F_r(s)=s\{r(4-s)+3\}-r<0$ である．
最後に式~\eqref{zu16:pair} と $s\ge1$ から
\[
 y<\frac{r+2}{r^2}S,\qquad
 E<\left(4+\frac4r-s\right)\frac Sr
\]
となる．$s\ge10$ なら右辺は $-2$ より小さいので，$E\ge-2$ に反する．
\end{proof}

\begin{lemma}[小さい $S$ の完全な処理]
\label{zu16:small-S}
$S\le3$ の場合，接続可能性を残す中央模型は
\[
 (r,s,S,T)=(1,1,2,6),\ (1,4,3,12),\
 (2,2,3,6),\ (2,4,3,6),\ (2,6,3,6)
\]
の五つに限られる．
\end{lemma}

\begin{proof}
$S>r$ と式~\eqref{zu16:pair} から，
$(r,S)=(1,2)$ では $(s,y)=(1,3)$，
$(r,S)=(1,3)$ では $(s,y)=(4,6)$ だけとなる．
$(r,S)=(2,3)$ では $y=2$ が強制され，$s$ は任意の正の偶数である．
この共鳴族では $E=5-3s/2$，$D^2_3=s-2$ である．

後者で $s\ge8$ と仮定する．両葉で ray $2=B$ は node 条件に反する．
ray $2=C$ の場合は $u=th=1$ が必要だが，係数相殺が非零の $d$ を要求し，
$c=0$ となる．すると他方の中心への entry が零になり，
葉 $3$ では中央 row の $T=6$，葉 $4$ では $r=2$ と両立しない．
従って両葉で ray $2=D$ である．node 条件は $2\le v\le4$ を与え，
両葉の parameter $h$ は三以下となる．
葉 $3$ では $c_3=s-h_3-3>0$，
葉 $4$ では $c_4=E-h_4-1<0$ なので，両方で $d=0$ である．
中央の正の entries $T,r$ を打ち消すには，他方の中心の ray が $B$，
他の葉の ray が $C$ でなければならない．
従って $D^3_4=h_3>0$，$D^4_3=h_4>0$ となり，相殺できない．
よって共鳴族では $s=2,4,6$ だけを調べればよい．
\end{proof}

以下，$S\ge4$ とする．葉 $3$ では低辺 $23$ の node 条件により
$p^3_{20}\ge S-4\ge0$ であり，高低 pair product が正なので実際は正である．
従ってその ray $2$ は $D$ である．葉 $3$ の parameters を $h,t,c,d$ と書くと
\begin{equation}\label{zu16:leaf3}
 v=t(h+1),\quad c=s-h-3,\quad
 \alpha=v-S\in[-4,1]\cap\ZZ,\quad
 \beta=p^3_{21}+T\in[-4,1]\cap\ZZ .
\end{equation}
第三の node 和 $\gamma=-2-\alpha-\beta$ も $[-4,1]\cap\ZZ$ に入る．

\begin{lemma}[非 $B$ 葉の有限分岐]
\label{zu16:nonB-split}
葉 $4$ の ray $2$ を $C$ または $D$ として表すと，次のいずれかが成立する．
\begin{enumerate}
\item $-2\le E\le7$．
\item $r=1$，葉 $4$ の ray $3$ が $B$ であり，
\[
 D^4_3=E+\kappa,\qquad \kappa\in\{2,5,8,10\}.
\]
\end{enumerate}
\end{lemma}

\begin{proof}
葉 $4$ の parameters を $H,\tau,c_4,d_4$ とする．
ray $2=C$ なら，node 条件は $1\le u_4=\tau H\le3-r$ を与える．
$c_4\ne0$ なら $d_4=0$，係数相殺より $E=-H$ なので第一の場合となる．
$c_4=0$ なら，残る二低 rays の割当に応じて
$(B,D)=(0,3)$ または $(3,0)$ である．
前者では $r=u_4=H=\tau=1$ であり，$D^4_3=-2$，$p^4_{31}=2$ である．
低辺 $34$ の puncture $1$ より葉 $3$ の ray $4$ は $C$，
$c\ge th>0$ となる．従って $d=0$ であり，
低辺 $23$ は $T\le v+2\le S+3$ を要求する．
一方 $r=1$ の正枝条件は $T>2S$ なので，$S\ge4$ に反する．
後者では $r=1$，$(H,\tau)=(1,1),(1,2),(2,1)$ であり，
\[
 D^4_3=E+\tau(H^2+H+1)+H+1
\]
から $\kappa=5,8,10$ を得る．

ray $2=D$ なら $2\le v_4=\tau(H+1)\le r+2\le6$ である．
$d_4\ne0$ なら $c_4=0$，係数相殺から $E=H+1\le6$ となる．
$d_4=0$，$(B,C)=(0,3)$ の場合は $D^4_3=H$ である．
葉 $3$ の ray $4$ が $C$ なら係数相殺により $d=h+H>0$，従って $c=0$ である．
ray $4=B$ でも $p^3_{42}=0$ なので，いずれも低辺 $34$ の puncture $2$ から
$-4\le c_4=E-H-1\le1$，すなわち $H-3\le E\le H+2$ を得る．
$1\le H\le5$ だから第一の場合である．
最後に $(B,C)=(3,0)$ の場合は $p^4_{20}=0$ なので $r=1$ であり，
$(H,\tau)=(1,1),(2,1)$ に限られる．
\[
 D^4_3=E+\tau(H^2+H+1)-H
\]
から $\kappa=2,5$ を得る．
\end{proof}

\begin{lemma}[自由な係数を残す分岐の有限化]
\label{zu16:free-E}
補題~\ref{zu16:nonB-split} の第二の場合には，次の有限集合だけを調べればよい．
\begin{enumerate}
\item 葉 $3$ の ray $4=C$：
$th\le4$，$-4\le c\le1$，$s=h+3+c$，
$S=t(h+1)-\alpha$，$-4\le\alpha\le1$．
\item 葉 $3$ の ray $4=B$：
$s=h+3$，$1\le h\le12$，$1\le t\le6$，
$S=t(h+1)-\alpha$ であり，node 和 $\alpha,\beta,\gamma$ と $\kappa$ は
\begin{equation}\label{zu16:free-E-identity}
 t(h^2+7h+3)
 =\alpha(h^2+5h+7)+(h+4)(1+\beta+\kappa)-2h-4
\end{equation}
を満たす．
\end{enumerate}
各場合で $y$ は式~\eqref{zu16:pair} により一意に決まり，$T=2y$ である．
\end{lemma}

\begin{proof}
葉 $4$ の ray $3=B$ の row は，高 ray $1$ への entry が $-1$，
他の低 rays への entries が零である．
従って葉 $3$ の ray $4=C$ の場合，低辺 $34$ の punctures $0,2$ が
$th\le4$，$-4\le c\le1$ を与える．

ray $4=B$ の場合，葉 $3$ の ray $1=C$ なので
$d=p^3_{21}=-T+\beta<0$，従って $c=0$，$s=h+3$ である．
低辺 $34$ の係数相殺は
\[
 T=t(h^2+h+1)+1+\beta+E+\kappa
\]
となる．これに $E=3y-(h+4)S$，$T=2y$，
$S=t(h+1)-\alpha$ と式~\eqref{zu16:pair} を代入して
式~\eqref{zu16:free-E-identity} を得る．

$A=h^2+7h+3$ と置く．右辺は
$h^2+15h+51$ 以下であり，$7A-(h^2+15h+51)=6h^2+34h-30>0$
なので $t\le6$ である．
$\alpha\le0$ なら右辺は $10h+44$ 以下なので
$h^2-3h-41\le0$，従って $h\le8$ である．
$\alpha=1,t\ge2$ なら $h^2-h-45\le0$，従って $h\le7$ である．
残る $\alpha=t=1$ では，$K=1+\beta+\kappa$ と置くと
$(K-4)h+4K=0$ である．正整数解は $K=2,h=4$ または $K=3,h=12$
だけなので，すべての場合に $h\le12$ となる．
\end{proof}

\begin{proposition}[非 $B$ 表現の完全な有限候補集合]
\label{zu16:nonB-finite}
葉 $4$ の ray $2$ を $C$ または $D$ として表せる場合には，
中央模型は証明から定まる三十一個の候補のいずれかに属する．
これらの中央模型に低三角形 $234$ の全整数葉を接続すると，
次の十模型，計五十二組だけが残る：
\[
\begin{array}{rrrr|r}
 r&s&S&T&\#\text{低三角形}\\\hline
 1&4&3&12&21\\
 1&7&4&18&3\\
 1&4&8&36&4\\
 1&15&12&64&3\\
 2&4&3&6&4\\
 2&4&12&30&5\\
 2&5&14&36&1\\
 4&4&18&30&1\\
 4&4&35&60&8\\
 4&4&52&90&2
\end{array}
\]
表の復元数には ray $2=B$ と表せる葉も含めており，
同じ標識付き行列の組は一度だけ数えている．
\end{proposition}

\begin{proof}
補題~\ref{zu16:small-S} の五模型を取る．$S\ge4$ の有界係数分岐では，
$r=1,\ldots,4$，$s=1,\ldots,9$，$E=-2,\ldots,7$ を取り，
式~\eqref{zu16:E} から $S$，式~\eqref{zu16:pair} から $y$ を復元する．
それぞれの整数性と元の正枝条件を課す．自由係数分岐は
補題~\ref{zu16:free-E} の有限集合を用いる．
元の row sums，pair products，奇数分母，正の weights を再確認すると，
三十一個の異なる中央模型を得る．その生成 witness は六十六個である．

各中央模型について，一高分母の完全標準形の約数復元と ray $2=B$ の例外復元を
実行し，低辺 $23,24,34$ の三つずつの node 条件と係数相殺を再確認する．
この固定中央模型の復元は，node 区間と正約数に由来する有限手順である．
\nolinkurl{verify_two_leaf_unit_finiteness_v16.py} は全有限 trials，
中央模型と葉の完全行列，生成 witness を
\nolinkurl{two_leaf_unit_finiteness_certificate_v16.json} に保存する．
この計算に中央・葉・中心間の分母の任意の上限はない．
本命題が扱っていない ray $2=B$ のみの表現は，別に処理する必要がある．
\end{proof}
